%% Đề 41 -- de-tuyen-sinh-lop-10-thpt-mon-toan-nam-2013-2014-so-gddt-lang-son.pdf
\deheader{ĐỀ SỐ 41}%
  {Kỳ thi Tuyển sinh vào lớp 10 THPT — Năm học 2013--2014}%
  {Sở GD\&ĐT Lạng Sơn | Chỉ phần Đại số}
\addcontentsline{toc}{section}{Đề 41 — Tuyển sinh vào lớp 10 THPT — Năm học 2013--2014}

\begin{de_cau}{Câu 1}{2,0 điểm}
\begin{enumerate}[label=\alph*)]
  \item Tính giá trị của biểu thức: $A = \sqrt{75} - \sqrt{48} + \sqrt{3}$.
  \item Rút gọn biểu thức: $B = \dfrac{x-4}{\sqrt{x}+2}$ với $x \ge 0$, $x \neq 4$.
\end{enumerate}
\end{de_cau}

\begin{loi_giai_cau}
\textbf{a)} Ta có:
\begin{align*}
A &= \sqrt{25\cdot 3} - \sqrt{16\cdot 3} + \sqrt{3} \\
  &= 5\sqrt{3} - 4\sqrt{3} + \sqrt{3} \\
  &= (5 - 4 + 1)\sqrt{3} = \boxed{2\sqrt{3}}.
\end{align*}

\textbf{b)} Với $x \ge 0, x \neq 4$, ta nhận thấy $x - 4 = (\sqrt{x})^2 - 2^2 = (\sqrt{x}-2)(\sqrt{x}+2)$. Do đó:
\[B = \dfrac{(\sqrt{x}-2)(\sqrt{x}+2)}{\sqrt{x}+2} = \boxed{\sqrt{x}-2}.\]
\end{loi_giai_cau}

\begin{de_cau}{Câu 2}{2,0 điểm}
\begin{enumerate}[label=\alph*)]
  \item Giải phương trình: $x^2 - 2x - 3 = 0$.
  \item Giải hệ phương trình:
  \[\begin{cases} 2x + y = 5 \\ x - y = 1 \end{cases}.\]
\end{enumerate}
\end{de_cau}

\begin{loi_giai_cau}
\textbf{a)} Ta có $a = 1, b = -2, c = -3$.

Nhận thấy $a - b + c = 1 - (-2) + (-3) = 0$, nên phương trình có hai nghiệm phân biệt:
\[x_1 = -1, \quad x_2 = -\dfrac{c}{a} = 3.\]
\textit{(Hoặc phân tích thành nhân tử: $(x - 3)(x + 1) = 0 \Leftrightarrow x = 3$ hoặc $x = -1$).}

Vậy tập nghiệm của phương trình là $S = \{-1; 3\}$.

\textbf{b)} Cộng từng vế hai phương trình của hệ:
\[(2x + y) + (x - y) = 5 + 1 \Leftrightarrow 3x = 6 \Leftrightarrow x = 2.\]
Thay $x = 2$ vào phương trình thứ hai $x - y = 1$, ta có:
\[2 - y = 1 \Leftrightarrow y = 1.\]
Vậy hệ phương trình có nghiệm duy nhất $\boxed{(x; y) = (2; 1)}$.
\end{loi_giai_cau}

\begin{de_cau}{Câu 3}{2,0 điểm}
Cho phương trình bậc hai: $x^2 - 2mx + m^2 - 1 = 0$ ($m$ là tham số).
\begin{enumerate}[label=\alph*)]
  \item Chứng minh rằng phương trình luôn có hai nghiệm phân biệt với mọi giá trị của $m$.
  \item Tìm $m$ để hai nghiệm $x_1, x_2$ thỏa mãn hệ thức: $x_1^2 + x_2^2 = 10$.
\end{enumerate}
\end{de_cau}

\begin{loi_giai_cau}
\textbf{a)} Phương trình có các hệ số: $a = 1 \neq 0, b' = -m, c = m^2 - 1$.

Biệt thức thu gọn:
\[\Delta' = (b')^2 - ac = (-m)^2 - 1\cdot(m^2 - 1) = m^2 - m^2 + 1 = 1.\]
Vì $\Delta' = 1 > 0$ với mọi $m \in \mathbb{R}$, nên phương trình luôn có hai nghiệm phân biệt $x_1, x_2$ với mọi $m$.

\textbf{b)} Theo định lý Vi-ét, ta có:
\[\begin{cases} x_1 + x_2 = 2m \\ x_1 x_2 = m^2 - 1 \end{cases}\]
Biến đổi hệ thức đã cho:
\begin{align*}
x_1^2 + x_2^2 = 10 &\Leftrightarrow (x_1 + x_2)^2 - 2x_1 x_2 = 10 \\
&\Leftrightarrow (2m)^2 - 2(m^2 - 1) = 10 \\
&\Leftrightarrow 4m^2 - 2m^2 + 2 = 10 \\
&\Leftrightarrow 2m^2 = 8 \\
&\Leftrightarrow m^2 = 4 \Leftrightarrow \boxed{m = \pm 2}.
\end{align*}
Vậy $m \in \{2; -2\}$ là các giá trị cần tìm.
\end{loi_giai_cau}

\begin{de_cau}{Câu 4}{2,0 điểm}
Một ca nô xuôi dòng từ bến A đến bến B mất 2 giờ, rồi ngược dòng từ B về A mất 3 giờ. Biết vận tốc của dòng nước là $5\text{ km/h}$. Tính quãng đường sông từ bến A đến bến B (coi vận tốc riêng của ca nô không đổi).
\end{de_cau}

\begin{loi_giai_cau}
Gọi vận tốc riêng của ca nô (khi nước yên lặng) là $v$ (km/h, điều kiện: $v > 5$).

Khi đó:
\begin{itemize}
  \item Vận tốc ca nô khi xuôi dòng là: $v + 5$ (km/h).
  \item Vận tốc ca nô khi ngược dòng là: $v - 5$ (km/h).
\end{itemize}
Quãng đường xuôi dòng từ A đến B là: $S = 2(v + 5)$ (km).

Quãng đường ngược dòng từ B về A là: $S = 3(v - 5)$ (km).

Vì quãng đường xuôi dòng và ngược dòng là như nhau nên ta có phương trình:
\begin{align*}
2(v + 5) = 3(v - 5) &\Leftrightarrow 2v + 10 = 3v - 15 \\
                    &\Leftrightarrow 3v - 2v = 10 + 15 \\
                    &\Leftrightarrow v = 25 \text{ (thỏa mãn điều kiện } v > 5).
\end{align*}
Vận tốc riêng của ca nô là $25\text{ km/h}$.

Độ dài quãng đường sông AB là:
\[S = 2 \times (25 + 5) = 2 \times 30 = \boxed{60\text{ km}}.\]
\end{loi_giai_cau}

\begin{de_cau}{Câu 5}{0,5 điểm}
Cho hai số thực dương $a, b > 0$ thỏa mãn điều kiện $a + b = 1$. Tìm giá trị nhỏ nhất của biểu thức:
\[T = \dfrac{1}{a} + \dfrac{4}{b}.\]
\end{de_cau}

\begin{loi_giai_cau}
\textbf{Cách 1 (Áp dụng BĐT AM-GM):}
Vì $a + b = 1$ và $a, b > 0$, nhân $T$ với $(a + b)$:
\begin{align*}
T &= \left(\dfrac{1}{a} + \dfrac{4}{b}\right)(a + b) \\
  &= 1 + \dfrac{b}{a} + \dfrac{4a}{b} + 4 \\
  &= 5 + \left(\dfrac{b}{a} + \dfrac{4a}{b}\right).
\end{align*}
Áp dụng bất đẳng thức AM-GM cho hai số dương $\dfrac{b}{a}$ và $\dfrac{4a}{b}$:
\[\dfrac{b}{a} + \dfrac{4a}{b} \ge 2\sqrt{\dfrac{b}{a}\cdot\dfrac{4a}{b}} = 2\sqrt{4} = 4.\]
Suy ra:
\[T \ge 5 + 4 = 9.\]
Đẳng thức xảy ra khi và chỉ khi:
\[\dfrac{b}{a} = \dfrac{4a}{b} \Leftrightarrow b^2 = 4a^2 \Leftrightarrow b = 2a \text{ (do } a, b > 0).\]
Kết hợp với $a + b = 1$, ta có $a + 2a = 1 \Leftrightarrow a = \dfrac{1}{3} \Rightarrow b = \dfrac{2}{3}$ (thỏa mãn $a,b > 0$).

\textbf{Cách 2 (Bất đẳng thức Cauchy-Schwarz dạng Engel):}
Áp dụng BĐT Cauchy-Schwarz dạng Engel cho hai phân thức:
\[T = \dfrac{1^2}{a} + \dfrac{2^2}{b} \ge \dfrac{(1 + 2)^2}{a + b} = \dfrac{3^2}{1} = 9.\]
Đẳng thức xảy ra khi $\dfrac{1}{a} = \dfrac{2}{b} \Leftrightarrow b = 2a$, kết hợp $a + b = 1 \Rightarrow a = \dfrac{1}{3}, b = \dfrac{2}{3}$.

Vậy giá trị nhỏ nhất của biểu thức là $\boxed{T_{\min} = 9}$, đạt được khi $a = \dfrac{1}{3}$ và $b = \dfrac{2}{3}$.
\end{loi_giai_cau}
